\par\Needspace{6\baselineskip}\medskip
\hypertarget{native-155-article111}{}\pdfbookmark[2]{111.}{article-111}
\begin{center}111.\end{center}
Supponamus secundo, duos locos dari per distantias suas a tribus planis in Sole sub angulis rectis se secantibus; designemus has distantias pro loco primo per $x$, $y$, $z$, pro secundo per $x'$, $y'$, $z'$, supponamusque planum tertium esse ipsam eclipticam, plani primi et secundi autem polos positivos in longitudine $N$ et $90^0+N$ sitos esse. Ita erit per art. 53, duobus radiis vectoribus
\SourceLeaf{156}per $r$, $r'$ designatis,
\[
\begin{aligned}
x&=r\cos u\cos(N-\AscendingNode)+r\sin u\sin(N-\AscendingNode)\cos i\\
y&=r\sin u\cos(N-\AscendingNode)\cos i-r\cos u\sin(N-\AscendingNode)\\
z&=r\sin u\sin i\\
x'&=r'\cos u'\cos(N-\AscendingNode)+r'\sin u'\sin(N-\AscendingNode)\cos i\\
y'&=r'\sin u'\cos(N-\AscendingNode)\cos i-r'\cos u'\sin(N-\AscendingNode)\\
z'&=r'\sin u'\sin i.
\end{aligned}
\]
Hinc sequitur
\[
\begin{aligned}
zy'-yz'&=rr'\sin(u'-u)\sin(N-\AscendingNode)\sin i\\
xz'-zx'&=rr'\sin(u'-u)\cos(N-\AscendingNode)\sin i\\
xy'-yx'&=rr'\sin(u'-u)\cos i.
\end{aligned}
\]
E combinatione formulae primae cum secunda habebitur $N-\AscendingNode$ atque $rr'\sin(u'-u)\sin i$, hinc et ex formula tertia prodibit $i$ atque $rr'\sin(u'-u)$.

Quatenus locus, cui coordinatae $x'$, $y'$, $z'$ respondent, tempore posterior supponitur, $u'$ maior quam $u$ fieri debet: quodsi itaque insuper constat, utrum angulus inter locum primum et secundum circa Solem descriptus duobus rectis minor an maior sit, $rr'\sin(u'-u)\sin i$ atque $rr'\sin(u'-u)$ esse debent quantitates positivae in casu primo, negativae in secundo: tunc itaque $N-\AscendingNode$ sine ambiguitate determinatur, simulque ex signo quantitatis $xy'-yx'$ deciditur, utrum motus directus sit, an retrogradus. Vice versa, si de motus directione constat, e signo quantitatis $xy'-yx'$ decidere licebit, utrum $u'-u$ minor an maior quam $180^0$ accipiendus sit. Sin vero tum motus directio, tum indoles anguli circa Solem descripti plane incognitae sunt, manifestum est, inter nodum ascendentem ac descendentem distinguere non licere.

Ceterum facile perspicitur, sicuti $\cos i$ est cosinus inclinationis plani orbitae versus planum tertium, ita $\sin(N-\AscendingNode)\sin i$, $\cos(N-\AscendingNode)\sin i$ esse resp. cosinus inclinationum plani orbitae versus planum primum et secundum; nec non exprimere $rr'\sin(u'-u)$ duplam aream trianguli inter duos radios vectores inclusi, atque $zy'-yz'$, $xz'-zx'$, $xy'-yx'$ duplam aream proiectionum eiusdem trianguli ad singula plana.

Denique patet, planum tertium pro ecliptica quodvis aliud planum esse posse, si modo omnes magnitudines per relationes suas ad eclipticam definitae perinde ad planum tertium, quidquid sit, referantur.
